
%%%%%%%%%%%%%%%%%%%%%%%%%%%%%%%%%%%%%%%%%%%%%%%%%%%%%%%%%%%%%%%%%%%%%%%%%%%%%%
%% OpmaakTikZ.tex
%% FVD Graphics Library — Ximera-proof basisversie
%%
%% Belangrijk:
%% - definieert GEEN nieuwe kleuren
%% - definieert GEEN nieuwe tikzset-stijlen
%% - gebruikt alleen bestaande stijlen uit Opmaak.tex:
%%     fvd axis
%%     fvd tick
%%     fvd guide
%%     fvd label
%%     fvd panel
%%
%% Deze versie vermijdt:
%% - externe afbeeldingen
%% - emoji
%% - background layers
%% - complexe pgfkeys
%% - berekeningen in tikz-coördinaten
%%%%%%%%%%%%%%%%%%%%%%%%%%%%%%%%%%%%%%%%%%%%%%%%%%%%%%%%%%%%%%%%%%%%%%%%%%%%%%

% ============================================================
% VEILIGHEID: VOORKOM DUBBEL INLADEN
% ============================================================

\ifdefined\fvdTikZLibraryLoaded
  \expandafter\endinput
\fi

\global\let\fvdTikZLibraryLoaded\relax


% ============================================================
% BASIS: GETALLENLIJN
% ============================================================
%
% Gebruik binnen een tikzpicture:
%
%   \FVDNumberLine{-6}{6}
%
% ============================================================

\newcommand{\FVDNumberLine}[2]{%
  \draw[fvd axis] (#1-0.6,0) -- (#2+0.6,0);

  \foreach \x in {#1,...,#2}{%
    \draw[fvd tick] (\x,0.13) -- (\x,-0.13);

    \ifnum\x=0
      \node[
        below=3pt,
        text=fvdmagenta,
        font=\small\bfseries
      ] at (\x,0) {0};
    \else
      \node[
        below=3pt,
        text=fvddark,
        font=\small
      ] at (\x,0) {\x};
    \fi
  }%
}


% ============================================================
% PUNTEN OP DE GETALLENLIJN
% ============================================================

\newcommand{\FVDPoint}[1]{%
  \fill[fvdmagenta] (#1,0) circle (2.2pt);
}

\newcommand{\FVDPointCyan}[1]{%
  \fill[fvdcyan] (#1,0) circle (2.2pt);
}

\newcommand{\FVDLabeledPoint}[2]{%
  \fill[fvdmagenta] (#1,0) circle (2.2pt);

  \node[
    fvd label,
    above=7pt
  ] at (#1,0) {#2};
}


% ============================================================
% HULPLIJNEN
% ============================================================

\newcommand{\FVDGuideUp}[2]{%
  \draw[fvd guide] (#1,0.08) -- (#1,#2);
}


% ============================================================
% EENVOUDIGE BEWEGING
% ============================================================
%
% Gebruik:
%
%   \FVDMoveLeft{5}{2}{3 stappen naar links}
%   \FVDMoveRight{2}{7}{5 stappen naar rechts}
%
% De macro's gebruiken een vaste booghoogte.
% Daardoor zijn ze betrouwbaarder in TeX4ht dan berekende control points.
%
% ============================================================

\newcommand{\FVDMoveLeft}[3]{%
  \draw[
    ->,
    line width=1.4pt,
    color=fvdcyan
  ]
    (#1,0.52)
    to[bend left=28]
    (#2,0.52);

  \node[
    fill=white,
    draw=fvdcyan!40,
    text=fvddark,
    rounded corners=2pt,
    inner xsep=5pt,
    inner ysep=2pt,
    font=\small
  ] at (#2+1.5,1.15) {#3};

  \FVDGuideUp{#1}{0.50}
  \FVDGuideUp{#2}{0.50}
}

\newcommand{\FVDMoveRight}[3]{%
  \draw[
    ->,
    line width=1.4pt,
    color=fvdcyan
  ]
    (#1,0.52)
    to[bend left=28]
    (#2,0.52);

  \node[
    fill=white,
    draw=fvdcyan!40,
    text=fvddark,
    rounded corners=2pt,
    inner xsep=5pt,
    inner ysep=2pt,
    font=\small
  ] at (#1+2.0,1.15) {#3};

  \FVDGuideUp{#1}{0.50}
  \FVDGuideUp{#2}{0.50}
}


% ============================================================
% PIXELPERSONAGE
% ============================================================
%
% Gebruik binnen een scope:
%
% \begin{scope}[shift={(5,0.75)},scale=0.06]
%   \FVDPixelPlayer
% \end{scope}
%
% ============================================================

\newcommand{\FVDPixelPlayer}{%
  % hoofd
  \fill[fvddark] (0,4) rectangle (4,8);

  % ogen
  \fill[fvdcyan] (0.8,5.7) rectangle (1.7,6.6);
  \fill[fvdcyan] (2.8,5.7) rectangle (3.7,6.6);

  \fill[white] (1.05,5.95) rectangle (1.45,6.35);
  \fill[white] (3.05,5.95) rectangle (3.45,6.35);

  % lichaam
  \fill[fvdmagenta] (1,1.5) rectangle (3.5,4);

  % armen
  \fill[fvdcyan] (0,2.3) rectangle (1,3.8);
  \fill[fvdcyan] (3.5,2.3) rectangle (4.5,3.8);

  % benen
  \fill[fvddark] (1,0) rectangle (1.9,1.5);
  \fill[fvddark] (2.6,0) rectangle (3.5,1.5);
}


% ============================================================
% PIXELCONTROLLER
% ============================================================
%
% Gebruik binnen een scope:
%
% \begin{scope}[shift={(0,0)},scale=0.06]
%   \FVDPixelController
% \end{scope}
%
% ============================================================

\newcommand{\FVDPixelController}{%
  % buitenvorm
  \fill[fvddark] (1,2) rectangle (11,7);
  \fill[fvddark] (0,3) rectangle (12,6);
  \fill[fvddark] (2,1) rectangle (4,3);
  \fill[fvddark] (8,1) rectangle (10,3);

  % binnenvlak
  \fill[white] (1.5,2.5) rectangle (10.5,6.5);
  \fill[fvdcyan] (1.5,6.0) rectangle (10.5,6.5);

  % d-pad
  \fill[fvddark] (2.2,4.1) rectangle (5.0,4.9);
  \fill[fvddark] (3.2,3.1) rectangle (4.0,5.9);

  % actieknoppen
  \fill[fvdmagenta] (8.0,4.8) circle (0.42);
  \fill[fvdmagenta] (9.2,3.7) circle (0.42);
  \fill[fvdcyan] (6.9,3.7) circle (0.42);
  \fill[fvdcyan] (8.0,2.8) circle (0.42);

  % middenknoppen
  \fill[fvddark] (5.2,4.15) rectangle (5.9,4.45);
  \fill[fvddark] (6.2,4.15) rectangle (6.9,4.45);
}


% ============================================================
% KANT-EN-KLARE FIGUUR: STARTPOSITIE
% ============================================================
%
% Gebruik:
%
%   \FVDStartPositionFigure{-6}{6}{5}
%
% ============================================================

\newcommand{\FVDStartPositionFigure}[3]{%
  \begin{center}
    \begin{tikzpicture}[
      x=0.85cm,
      y=1cm,
      line cap=round,
      line join=round
    ]


      \FVDNumberLine{#1}{#2}

      \draw[fvd guide]
        (#3,0.08)
        --
        (#3,0.70);

      \begin{scope}[shift={(#3-0.18,0.76)},scale=0.055]
        \FVDPixelPlayer
      \end{scope}

      \node[
        fvd label
      ] at (#3,1.42) {startpositie \(#3\)};

      \node[
        anchor=east,
        text=fvdmagenta,
        font=\small\bfseries
      ] at (-0.35,-1) {negatieve richting};

      \node[
        anchor=west,
        text=fvdcyan,
        font=\small\bfseries
      ] at (0.35,-1) {positieve richting};

    \end{tikzpicture}
  \end{center}
}


% ============================================================
% KANT-EN-KLARE FIGUREN: BEWEGING
% ============================================================
%
% Deze twee macro's zijn bewust apart gehouden.
% Zo vermijden we complexe berekeningen en conditionele TikZ-code.
%
% ============================================================

\newcommand{\FVDMovementLeftFigure}[5]{%
  \begin{center}
    \begin{tikzpicture}[
      x=0.85cm,
      y=1cm,
      line cap=round,
      line join=round
    ]

        (#1-0.8,-0.72)
        rectangle
        (#2+0.8,1.72);

      \FVDNumberLine{#1}{#2}

      \FVDPoint{#3}
      \FVDPointCyan{#4}
      \FVDMoveLeft{#3}{#4}{#5}

    \end{tikzpicture}
  \end{center}
}

\newcommand{\FVDMovementRightFigure}[5]{%
  \begin{center}
    \begin{tikzpicture}[
      x=0.85cm,
      y=1cm,
      line cap=round,
      line join=round
    ]

      \fill[fvdcyan!6]
        (#1-0.8,-0.72)
        rectangle
        (#2+0.8,1.72);

      \FVDNumberLine{#1}{#2}

      \FVDPoint{#3}
      \FVDPointCyan{#4}
      \FVDMoveRight{#3}{#4}{#5}

    \end{tikzpicture}
  \end{center}
}


% ============================================================
% EENVOUDIG ASSENSTELSEL
% ============================================================
%
% Gebruik binnen een tikzpicture:
%
%   \FVDCoordinatePlane{-5}{5}{-4}{4}
%
% ============================================================

\newcommand{\FVDCoordinatePlane}[4]{%
  \draw[
    step=1cm,
    color=fvdcyan!18,
    line width=0.35pt
  ] (#1,#3) grid (#2,#4);

  \draw[fvd axis]
    (#1,0)
    --
    (#2+0.35,0)
    node[
      below left=2pt,
      text=fvddark
    ] {\(x\)};

  \draw[fvd axis]
    (0,#3)
    --
    (0,#4+0.35)
    node[
      below left=2pt,
      text=fvddark
    ] {\(y\)};

  \foreach \x in {#1,...,#2}{%
    \ifnum\x=0
    \else
      \draw[fvd tick] (\x,0.10) -- (\x,-0.10);
      \node[
        below=2pt,
        text=fvddark,
        font=\scriptsize
      ] at (\x,0) {\x};
    \fi
  }

  \foreach \y in {#3,...,#4}{%
    \ifnum\y=0
    \else
      \draw[fvd tick] (0.10,\y) -- (-0.10,\y);
      \node[
        left=2pt,
        text=fvddark,
        font=\scriptsize
      ] at (0,\y) {\y};
    \fi
  }

  \node[
    below left=2pt,
    text=fvdmagenta,
    font=\scriptsize\bfseries
  ] at (0,0) {0};
}


% ============================================================
% PUNT IN EEN ASSENSTELSEL
% ============================================================

\newcommand{\FVDCoordinatePoint}[3]{%
  \draw[fvd guide] (#1,0) -- (#1,#2);
  \draw[fvd guide] (0,#2) -- (#1,#2);

  \fill[fvdmagenta] (#1,#2) circle (2.2pt);

  \node[
    fvd label,
    above right=3pt
  ] at (#1,#2) {\(#3\)};
}


%%%%%%%%%%%%%%%%%%%%%%%%%%%%%%%%%%%%%%%%%%%%%%%%%%%%%%%%%%%%%%%%%%%%%%%%%%%%%%
%% EINDE
%%%%%%%%%%%%%%%%%%%%%%%%%%%%%%%%%%%%%%%%%%%%%%%%%%%%%%%%%%%%%%%%%%%%%%%%%%%%%%
